\documentclass[11pt]{article}
\usepackage{varnernotes}

\newcommand{\R}{\mathbb{R}}

\begin{document}

\lecturenotes{10b}{Probability of Profit, Composite Contracts, and Delta Hedging}{Fall 2026}{October 29, 2026}{Jeffrey D. Varner}

\learningobjectives
  {Aggregate option legs}
  {Construct composite expiration payoff and profit functions with consistent signs, quantities, premiums, and contract multipliers.}
  {Compute strategy probabilities}
  {Map profitable terminal-price regions through a stated predictive distribution and distinguish profit from exercise-related events.}
  {Build a local hedge}
  {Sum leg deltas, select an offsetting stock position, and explain why delta neutrality requires monitoring and rebalancing.}

\section{A Strategy Is the Sum of Signed Legs}

Options are combined to reshape exposure rather than merely add contracts. Let a composite position $\mathcal C$ contain legs $i=1,\ldots,d$ on the same underlying. Leg $i$ has type $\tau_i\in\{c,p\}$, strike $K_i>0$ in USD/share, signed orientation $\theta_i\in\{-1,+1\}$ for short/long, count $n_i\in\mathbb N$, and opening premium $P_i\geq 0$ in USD/share. Define the expiration payoff per share by
\begin{equation}
  \label{eq:leg-payoff}
  h_i(S_T):=
  \begin{cases}
    (S_T-K_i)^+, & \tau_i=c,\\
    (K_i-S_T)^+, & \tau_i=p,
  \end{cases}
  \qquad x^+:=\max\{x,0\},
\end{equation}
where $S_T>0$ is the underlying price at common expiration $T$. With a multiplier $q>0$ shares/contract, terminal payoff and frictionless profit are
\begin{align}
  \label{eq:composite-payoff}
  H_{\mathcal C}(S_T)
  &:=q\sum_{i=1}^{d}\theta_i n_i h_i(S_T),\\
  \label{eq:composite-profit}
  \Pi_{\mathcal C}(S_T)
  &:=q\sum_{i=1}^{d}\theta_i n_i\bigl[h_i(S_T)-P_i\bigr].
\end{align}
Equation~\eqref{eq:composite-profit} treats a long premium as an initial cost and a short premium as an initial receipt. Commissions, bid--ask spread, financing, dividends on stock hedges, assignment fees, taxes, and slippage must be added for an economic P\&L.

\begin{remark}{Expiration P\&L versus mark-to-market P\&L}{pnl-horizon}
Before expiration, replace $h_i(S_T)$ in the closing value with an executable option mark $V_i(S_t,t)$ and include hedge cash flows. Comparing today's spot directly with an expiration breakeven does not determine current P\&L because remaining extrinsic value and close-out spreads still matter.
\end{remark}

\section{Straddles: Two Tails and Two Breakevens}

A one-contract long straddle buys a call and put with common strike $K$ and expiration. Let $D:=P_c+P_p>0$ be the total debit per share. Because $(S_T-K)^++(K-S_T)^+=|S_T-K|$, its profit is
\begin{equation}
  \label{eq:straddle-profit}
  \Pi_L(S_T)=q\bigl(|S_T-K|-D\bigr).
\end{equation}
The breakevens are $B_-:=K-D$ and $B_+:=K+D$, provided $B_->0$ for the lower boundary to be economically reachable. The long straddle profits when $S_T<B_-$ or $S_T>B_+$; its maximum loss is $qD$ at $S_T=K$, and its upside profit is unbounded in the idealized model. The short straddle has profit $-\Pi_L$; it collects at most $qD$ but has unbounded upside loss and large downside loss.

\begin{figure}[ht]
  \centering
  \begin{tikzpicture}[x=0.72cm,y=0.56cm,>=Latex,font=\sffamily\small]
    \draw[->,vnmuted] (0,2.2)--(10.3,2.2) node[right] {$S_T$};
    \draw[->,vnmuted] (5.0,0)--(5.0,6.0) node[above] {profit};
    \draw[vnlink,very thick] (0.7,5.35)--(5.0,0.8)--(9.4,5.45);
    \draw[vncarnelian,very thick,dashed] (0.7,-0.95)--(5.0,3.6)--(9.4,-1.05);
    \draw[vnmuted,dotted] (2.9,0)--(2.9,5.6);
    \draw[vnmuted,dotted] (7.1,0)--(7.1,5.6);
    \node[below] at (2.9,2.2) {$B_-$};
    \node[below] at (5.0,2.2) {$K$};
    \node[below] at (7.1,2.2) {$B_+$};
    \node[vnlink,above right] at (8.1,4.15) {long};
    \node[vncarnelian,below right] at (7.7,0.55) {short};
  \end{tikzpicture}
  \caption{Schematic expiration profit for long and short straddles. The long position buys both tails beyond $B_-$ and $B_+$; the short position sells those tails in exchange for the initial credit.}
  \label{fig:straddle-profit}
\end{figure}

Let $F(x):=\mathbb P(S_T\leq x)$ be the CDF under a clearly stated predictive model. If $F$ is continuous at both breakevens, then
\begin{equation}
  \label{eq:straddle-pop}
  \operatorname{POP}_{L}=F(B_-)+1-F(B_+),
  \qquad
  \operatorname{POP}_{S}=F(B_+)-F(B_-).
\end{equation}
The two sum to one only under matching frictionless profit functions and zero probability at the breakevens. With discrete distributions, use strict and non-strict inequalities consistently.

\notebooklink
  {L10b: AMD Earnings Straddle}
  {https://github.com/varnerlab/CHEME-5660-CourseRepository-Fall-2026/blob/main/lectures/week-10/L10b/CHEME-5660-L10b-Example-Straddle-AMD-Earnings-Fall-2026.ipynb}
  {Build long and short straddle profit diagrams, locate both breakevens, and estimate terminal profit probabilities from projected price distributions.}

\section{Vertical Put Credit Spreads: Bounded Tails}

Sell one put at $K_H$ for premium $P_H$ and buy one put at $K_L<K_H$ for premium $P_L$. Let the net credit be $C:=P_H-P_L>0$ and the width be $W:=K_H-K_L>0$. Per share,
\begin{equation}
  \label{eq:put-spread-profit}
  \pi_V(S_T)=C+(K_L-S_T)^+-(K_H-S_T)^+.
\end{equation}
The maximum profit is $C$ for $S_T\geq K_H$; the minimum profit is $C-W$ for $S_T\leq K_L$, so maximum \emph{loss magnitude} is $W-C$ when $0<C<W$. The breakeven is
\begin{equation}
  \label{eq:put-spread-breakeven}
  B:=K_H-C,
\end{equation}
and profit occurs for $S_T>B$. Under CDF $F$, $\operatorname{POP}_V=1-F(B)$ when the boundary has zero probability. This defined-risk structure caps both reward and downside, but early assignment, pin risk, and illiquid close-out prices still matter before expiration.

\begin{figure}[ht]
  \centering
  \begin{tikzpicture}[x=0.85cm,y=0.6cm,>=Latex,font=\sffamily\small]
    \draw[->,vnmuted] (0,2.7)--(9.6,2.7) node[right] {$S_T$};
    \draw[->,vnmuted] (0.5,0)--(0.5,5.6) node[above] {profit};
    \draw[vnlink,very thick] (0.8,1.0)--(3.4,1.0)--(6.6,4.25)--(9.1,4.25);
    \draw[vnmuted,dotted] (3.4,0.25)--(3.4,5.1);
    \draw[vnmuted,dotted] (5.1,0.25)--(5.1,5.1);
    \draw[vnmuted,dotted] (6.6,0.25)--(6.6,5.1);
    \node[below] at (3.4,2.7) {$K_L$};
    \node[below] at (5.1,2.7) {$B$};
    \node[below] at (6.6,2.7) {$K_H$};
    \node[vnlink,above] at (8.0,4.25) {max profit $C$};
    \node[vnlink,below] at (2.2,1.0) {min profit $C-W$};
  \end{tikzpicture}
  \caption{Expiration profit per share for a put credit vertical. The long lower-strike put truncates the short higher-strike put's downside.}
  \label{fig:vertical-profit}
\end{figure}

\notebooklink
  {L10b: Put Credit Vertical Spread and Delta}
  {https://github.com/varnerlab/CHEME-5660-CourseRepository-Fall-2026/blob/main/lectures/week-10/L10b/CHEME-5660-L10b-Example-Delta-VerticalSpread-Fall-2026.ipynb}
  {Combine two put legs, compute bounded profit and loss, estimate terminal probabilities, and aggregate the spread's delta.}

\section{Probability Requires a Measure and a Horizon}

The definition
\begin{equation}
  \label{eq:general-pop}
  \operatorname{POP}_{\mathcal C}^{\mathbb P}
  :=\mathbb P\!\left(\Pi_{\mathcal C}(S_T)>0\mid\mathcal I_t\right)
\end{equation}
is incomplete until the probability measure $\mathbb P$, information set $\mathcal I_t$, horizon, parameters, and cost convention are specified. A historical or forecast model attempts to describe real-world outcomes. The risk-neutral measure $\mathbb Q$ is constructed for arbitrage pricing; $\mathbb Q$-probabilities are not automatically real-world frequencies or personal profit odds.

Historical GBM, option-implied distributions, and simple normal ``expected move'' approximations answer different questions and embed different assumptions. Historical estimates are backward-looking and unstable. Implied distributions are affected by risk premia and require a full option surface, not a single volatility quote. A normal price approximation can assign mass to negative prices. Report probabilities as model-conditional estimates and stress the distribution, not as objective facts.

For exercise-related events, distinguish terminal ITM from American exercise. A call is terminally ITM when $S_T>K$ and a put when $S_T<K$. That event does not include early exercise, broker assignment rules, or exercise-by-exception thresholds. In a composite position, report leg-level events or an explicitly defined joint event; there is no universal strategy-level ``probability of exercise.''

\section{Delta Hedging a Composite Position}

Let $V_i(S,t)$ be the current per-share value of leg $i$, let $\Delta_i:=\partial V_i/\partial S$, and let $N_S\in\R$ be the number of underlying shares held. The marked portfolio value and net delta are
\begin{align}
  \label{eq:portfolio-value}
  \mathcal V(S,t)
  &:=N_SS+q\sum_{i=1}^{d}\theta_i n_iV_i(S,t),\\
  \label{eq:net-delta}
  \Delta_{\mathrm{net}}
  &:=\frac{\partial\mathcal V}{\partial S}
   =N_S+q\sum_{i=1}^{d}\theta_i n_i\Delta_i.
\end{align}
At the current state $(S,t)$, choose
\begin{equation}
  \label{eq:hedge-shares}
  N_S^*:=-q\sum_{i=1}^{d}\theta_i n_i\Delta_i
\end{equation}
to make $\Delta_{\mathrm{net}}=0$. If shares must be integral, round deliberately and report residual delta. For ten short calls with $q=100$ and long-call delta $0.60$, options delta is $10(-1)(100)(0.60)=-600$ shares, so holding $600$ shares offsets the initial first-order exposure.

\begin{theorem}{Delta neutrality is local}{local-hedge}
A stock holding selected by \eqnref{hedge-shares} neutralizes the first derivative at the current valuation state. Unless the options position has zero gamma and unchanged inputs, the same fixed $N_S^*$ will not keep $\Delta_{\mathrm{net}}=0$ for all future $S$ and $t$. Dynamic delta hedging therefore requires recalculation and trading.
\end{theorem}

Delta hedging does not eliminate gamma, volatility, time, rate, jump, liquidity, funding, or model risk. Rebalancing more frequently reduces some discretization error but increases turnover and costs. Early exercise or assignment changes both the option inventory and stock position immediately, so the hedge and cash account must be updated as contract state changes.

\Needspace{0.38\textheight}
\keytakeaways
  {In this lecture, we built composite expiration P\&L, mapped profitable regions through explicit distributions, and constructed a local stock hedge from leg deltas.}
  {Profit regions come from the full strategy}
  {Add signed leg payoffs and opening cash flows first; then solve $\Pi_{\mathcal C}(S_T)>0$ for the relevant terminal-price region.}
  {POP is model-conditional}
  {A probability needs a measure, horizon, information set, and cost definition; risk-neutral pricing probabilities are not automatically real-world profit odds.}
  {Delta neutrality is temporary}
  {The hedge cancels first-order spot exposure at the current state, while gamma, time, other factors, costs, and assignment make rebalancing necessary.}
  {Next, apply payoff, risk, and assignment logic to covered calls and cash-secured puts.}

\begin{readinglist}
  \item Hayne E. Leland, \href{https://doi.org/10.1111/j.1540-6261.1985.tb02383.x}{``Option Pricing and Replication with Transactions Costs,''} \emph{Journal of Finance} 40(5), 1283--1301 (1985).
  \item Mark Rubinstein and Eric Reiner, \emph{Exotic Options}, working-paper and practitioner treatments of static payoff decomposition (1991).
  \item John C. Hull, \emph{Options, Futures, and Other Derivatives}, 9th ed., Pearson (2015), chapters on trading strategies and the Greeks.
\end{readinglist}

\vfill
{\sffamily\footnotesize\color{vnmuted}\textbf{Educational use and risk notice.} These notes are for instruction only and do not constitute investment advice. Probability estimates and hedges are model-dependent, and options trading involves substantial risk.}

\end{document}
